% Results (main paper). Numbers: tab1_controlled, tab2_published, tab3_ablation, tab4_coverage, tab5_reasoning,
% tab6_editing; appendix pointers: tab:backbones, tab:multitarget, tab:tree_ablation, tab:reason_prompts.
% All numbers are test-split results.
\subsection{Main results}
\label{sec:main_results}

\paragraph{Controlled comparison.}
Tab.~\ref{tab:controlled} compares all methods under the one-mask rule with the same $4$B backbone. \ours{} is the
best row on $13$ of $14$ columns. The gains are largest on parts, where existing systems struggle most: $+14.5$
gIoU on PACO, $+13.9$ on GSEval, $+7.1$ on Pascal-Part-116, $+5.6$ on PartImageNet and $+2.3$ on PARTONOMY over the
best baseline, all significant ($p<10^{-4}$). On InstructPart the difference ($+0.8$) is not significant, and we
count it as a tie. The table also shows why parts are hard for a box pipeline. Box\,+\,SAM\,3, the strongest
baseline, reaches $76.4$ on PACO objects but $33.0$ on PACO parts: the MLLM finds the right place, but a box around
a part often makes the segmenter return the whole object (Fig.~\ref{fig:showcase}). \ours{} avoids this because it
chooses among masks that already have the right extent.

These gains do not come at the cost of whole objects. \ours{} improves gRefCOCO by $+8.0$ and $+7.9$ gIoU and PACO
objects by $+5.4$; the only column where it is behind is the full GSEval set ($-2.3$). The clearest change in
behaviour is rejection. The MLLM baselines almost always return a mask and reject at most $32.4\%$ of the
expressions with no target, while \ours{} rejects $94.0\%$ and $90.9\%$. SAM\,3 also rejects well, but only by
giving up grounding accuracy. The same picture holds with an $8$B backbone, and a $32$B Box\,+\,SAM\,3 pipeline
still stays below our $4$B model on gRefCOCO and on rejection (Tab.~\ref{tab:backbones}). When several masks per
expression are allowed, \ours{} also has the highest gIoU and N-acc on gRefCOCO (Tab.~\ref{tab:multitarget}).

\paragraph{Comparison with published systems.}
Tab.~\ref{tab:published} places \ours{} next to published results. Among models of up to $4$B it is the best on
PARTONOMY ($51.0$; the next system reaches $41.6$ with a $13$B backbone), on Pascal-Part-116, where it is
in-domain, and on rejection ($89.6$ / $86.0$ N-acc). It is not the best everywhere: PANORAMA and SAMTok are ahead on
gRefCOCO, PANORAMA on GSEval, and OP-HRG on PartImageNet and InstructPart. PANORAMA and SAMTok train $0.9$B to
$1.3$B parameters on one to five million samples. \ours{} trains $64.2$M parameters, $14$ to $21$ times fewer, on
three datasets, and still reaches the same range.

\subsection{Analysis}
\label{sec:analysis}

\paragraph{What each component contributes.}
Tab.~\ref{tab:ablation} changes one design choice at a time. Starting from SAM\,3 ranking its own text instances
(a), letting the MLLM choose among query-free masks (b) adds $13.1$ gIoU on PARTONOMY and $19.6$ on PartImageNet,
but loses on InstructPart and Pascal-Part, because the grid alone often lacks a clean mask of a small part. Adding
the query-conditioned masks (b$\rightarrow$e) repairs this and improves every part benchmark, by $3$ to $12$
points. Scoring pooled region features without mask tokens (c) costs $17.6$ gIoU on GSEval parts and $24.6$ on the
full GSEval set against (e), so reading the candidates inside the MLLM is the most important choice. Free
against region-restricted attention, (d) against (e), is mixed: free attention is better on GSEval ($+6.4$ on
parts), restricted attention on gRefCOCO testA ($+3.8$), and the other columns are within one point; we keep the
variant selected on validation. The tree-relation terms mainly help rejection, by $4.1$ and $5.0$ N-acc
(Tab.~\ref{tab:tree_ablation}).

\paragraph{Where the remaining errors come from.}
Tab.~\ref{tab:coverage} separates coverage from selection. On PARTONOMY the best mask in the pool reaches $71.0$
gIoU, \ours{} selects masks worth $49.4$, and SAM\,3's own top instance is worth $42.0$. The pool is therefore not
the limit: the right mask is usually present, and the remaining errors are errors of choice. This leaves clear
room for better selection over the same candidates.

\subsection{Reasoning queries}
\label{sec:reasoning}

Reasoning queries describe the target by its use, such as ``which part can be used to pick up the knife?'', and do
not name it, so they give the phrase-based half of our pool little to work with. For these benchmarks we add
\emph{query rewriting}: the base Qwen3-VL-4B first names the target as a short noun phrase, and this phrase
prompts SAM\,3. The selection is unchanged and still reads the original request. The rule was fixed on the
validation split of ReasonSeg before any test run, and the rewriter is the model Box\,+\,SAM\,3 uses for its box.

On InstructPart, where the target is a part, \ours{} with rewriting exceeds Box\,+\,SAM\,3 by $+11.2$ and $+8.5$
gIoU ($p<10^{-4}$; Tab.~\ref{tab:reasoning}). On ReasonSeg, where the target is a whole object, the two are on par
($-0.9$ gIoU, $p=0.39$). Without rewriting, \ours{} matches Box\,+\,SAM\,3 on InstructPart but is $9.7$ gIoU behind
on ReasonSeg, which shows that the missing piece was the candidate pool and not the selection. These conclusions
do not depend on the prompt given to the baselines (Tab.~\ref{tab:reason_prompts}).

\subsection{Region-grounded editing}
\label{sec:editing}

A mask is often an intermediate result that an editor uses to decide which pixels to change.
Tab.~\ref{tab:editing} uses each method as the mask source for two editors. \ours{} returns a mask for every
request, while Box\,+\,SAM\,3 returns none for $1.3\%$ and $2.8\%$ of them. Its mask IoU is the highest on
ImgEdit-Bench ($67.5$) and below Box\,+\,SAM\,3 on GEdit-Bench ($53.6$ against $57.6$). In edit quality the mask
sources are close: \ours{} has the best judge score in three of four settings, and the change outside the target
region is similar except with FLUX Fill on GEdit-Bench, where it is higher for \ours{}. The differences are small
and the judge favours larger edit regions, so we read this table as showing that masks selected by \ours{} are as
usable for editing as those of the baselines, not that they give better edits.
